\documentclass[12pt,a4paper]{article}
\usepackage[utf8]{inputenc}
\usepackage[T5]{fontenc}
\usepackage{geometry}
\usepackage{setspace}
\usepackage{graphicx}
\usepackage{booktabs}
\usepackage{float}
\usepackage{amsmath}
\usepackage{tabularx}
\usepackage{longtable}
\usepackage{array}
\geometry{margin=1in}
\onehalfspacing
\setcounter{tocdepth}{2}
\renewcommand{\arraystretch}{1.15}

\title{Báo cáo dự án xử lý ngôn ngữ tự nhiên\\ Phân tích cảm xúc cho các bài đánh giá thương mại điện tử bằng kỹ thuật học máy.}
\author{Team SkyDom+}
\date{February 2026}

\begin{document}

\begin{titlepage}
\centering
\vspace*{1.0cm}

\includegraphics[width=0.62\textwidth]{p01_chart1.png}\par
\vspace{1.0cm}
{\LARGE Phân tích tình cảm cho các bài đánh giá thương mại điện tử\par}
\vspace{0.2cm}
{\LARGE sử dụng Kỹ thuật học máy.\par}
\vspace{0.8cm}
{\Large Báo cáo nghiên cứu trường hợp\par}
\vspace{0.8cm}
{\large Đội: Đội SkyDom+\par}
\vspace{0.3cm}
{\large Thành viên nhóm:\par}
\vspace{0.4cm}
{\large
\begin{tabular}{@{}>{\raggedright\arraybackslash}p{5.8cm} >{\raggedright\arraybackslash}p{2.7cm}@{}}
Nguyen Nhu Tien Tuan & 22BA13314 \\
Tran Hien Chuong & 22BA13056 \\
Do Gia Bao & 22BA13042 \\
Bui The Anh & 22BA13007 \\
Le Duc An & 22BA13002 \\
Pham Quang Tung & 22BA13320 \\
Khuong Bao Duc & 22BA13077 \\
\end{tabular}\par}
\vfill
{\large Ngày 25 tháng 2 năm 2026\par}
\end{titlepage}

\pagenumbering{roman}

\begin{abstract}
Báo cáo này trình bày một dự án NLP để xem xét hiểu biết về dữ liệu Thẻ quà tặng Amazon. Dự án có ba lớp NLP. Lớp đầu tiên là phân loại cảm xúc ở cấp độ câu với các tính năng văn bản nhận biết ngữ cảnh. Lớp thứ hai là trích xuất vấn đề đa nhãn để xác định các loại khiếu nại cụ thể. Lớp thứ ba là một phần mở rộng biến đổi để phân tích ngữ cảnh sâu sắc và trường hợp khó. Tác phẩm sử dụng 152.410 bài đánh giá thô, với nhãn tình cảm mạnh mẽ từ xếp hạng 1, 2, 4 và 5, cũng như bộ phân tích không chắc chắn từ xếp hạng 3. Biến thể tình cảm cổ điển được chọn sử dụng từ và ký tự TF-IDF với gắn thẻ phủ định. Trên dữ liệu thử nghiệm, nó đạt độ chính xác đối với lớp âm là 0,677, thu hồi đối với lớp âm là 0,932 và F2 đối với lớp âm là 0,867. Phần mở rộng máy biến áp đạt được độ chính xác tổng thể cao hơn và độ chính xác âm cao hơn, nhưng mức thu hồi âm thô thô ở ngưỡng cố định thấp hơn. Để trích xuất vấn đề, mô hình nhiều nhãn cấp độ từ được kiểm tra micro-F1 là 0,765 và mô hình từ cộng với ký tự đạt 0,763. Để cải thiện sự liên kết khóa học mà không cần thêm dữ liệu, dự án hiện bao gồm gói nâng cấp giáo trình mạnh mẽ với Naive Bayes, Perceptron, LinearSVM, mất log SGD, đường cơ sở SVD ngữ nghĩa và đường cơ sở MLP chuyển tiếp nguồn cấp dữ liệu. Trong gói đó, mức thu hồi thử nghiệm\_0 dao động từ 0,719 đến 0,941, với F2\_0 tốt nhất là 0,855 cho Hồi quy logistic cân bằng và mức thu hồi tốt nhất\_0 là 0,941 cho Complement Naive Bayes. Gói này cũng bổ sung thêm các mô hình ngôn ngữ unigram và bigram với khả năng làm mịn add-k, trong đó độ phức tạp của thử nghiệm bigram cải thiện từ 377,45 lên 362,72. Một ma trận phù hợp với khóa học được xây dựng từ các hiện vật thực tế báo cáo mức độ bao phủ giáo trình ước tính là 75,7\%.
\end{abstract}

\tableofcontents
\newpage
\pagenumbering{arabic}

\section{Giới thiệu}
Ngôn ngữ tự nhiên trong đánh giá của người dùng khá ồn ào, ngắn gọn và hỗn tạp. Một đánh giá có thể bao gồm khen ngợi, phàn nàn, tiếng lóng, lỗi đánh máy và các từ miền trong cùng một câu. Điều này làm cho việc đánh giá hiểu được vấn đề NLP tốt cho cả phương pháp cổ điển và hiện đại \cite{jurafsky2026,turney2002}.

Dự án này nghiên cứu văn bản đánh giá từ Thẻ quà tặng Amazon. Mục tiêu không chỉ là gán nhãn tích cực hay tiêu cực. Mục tiêu là xây dựng một hệ thống ngôn ngữ thực tế có thể hiểu được cảm xúc, đánh dấu các trường hợp không chắc chắn và trích xuất các thẻ vấn đề cụ thể cho công việc hỗ trợ.

Dự án được thiết kế như một hệ thống NLP đầu cuối. Nó bao gồm chuẩn hóa văn bản, mô hình hóa tính năng từ vựng và từ phụ, xử lý phủ định và tương phản theo ngữ cảnh, học có giám sát, kiểm soát độ không đảm bảo, trích xuất vấn đề đa nhãn và mở rộng dựa trên máy biến áp.

\section{Định nghĩa vấn đề NLP}
\subsection{Nhiệm vụ A: Hiểu biết về tình cảm}
Nhiệm vụ A là phân loại cảm xúc nhị phân từ văn bản đánh giá. Xếp hạng 1 và 2 được ánh xạ tới lớp phủ định. Xếp hạng 4 và 5 được ánh xạ tới lớp tích cực. Xếp hạng có giá trị 3 được sử dụng để phân tích độ không chắc chắn và không phải là một phần của khóa đào tạo được giám sát chặt chẽ \cite{turney2002,wang2012}.

Chính sách quyết định thực tế không phải là tính chính xác đơn giản. Trong phân loại hỗ trợ, việc thiếu một đánh giá xấu là rất rủi ro. Do đó, việc lựa chọn mô hình sử dụng quy tắc phủ định trước với thứ tự ưu tiên thu hồi cho lớp 0, sau đó là F2 cho lớp 0, sau đó là độ chính xác cho lớp 0 \cite{saito2015}.

\subsection{Nhiệm vụ B: Trích xuất vấn đề}
Nhiệm vụ B là trích xuất vấn đề đa nhãn. Một đánh giá có thể chứa nhiều vấn đề nên việc gắn thẻ một nhãn là không đủ. Bộ nhãn đầu ra có chín loại vấn đề và cho phép gán nhiều nhãn trên mỗi hàng \cite{zhang2014,tsoumakas2011}.

Nhiệm vụ này là ánh xạ NLP cốt lõi từ văn bản đánh giá phi cấu trúc đến các tín hiệu vấn đề có cấu trúc. Nó giúp các nhóm hỗ trợ chuyển từ tình cảm chung sang các hạng mục hành động.

\subsection{Nhiệm vụ C: Mở rộng bối cảnh sâu}
Nhiệm vụ C là tinh chỉnh máy biến áp dưới dạng phần mở rộng NLP. Phần mở rộng này nhằm mục đích kiểm tra mô hình hóa ngôn ngữ theo ngữ cảnh sâu sắc đối với các cụm từ ngắn và mẫu phủ định khó. Nó cũng cung cấp sự so sánh trực tiếp với các biểu diễn thưa thớt cổ điển \cite{vaswani2017,devlin2019}.

\section{Đặc điểm dữ liệu và ngôn ngữ}
\subsection{Nguồn dữ liệu và quy mô}
Tệp nguồn chứa 152.410 đánh giá với các trường văn bản và xếp hạng. Bảng \ref{tab:nlp_rating_distribution} hiển thị phân bổ xếp hạng đầy đủ.

\begin{table}[H]
\centering
\caption{Phân phối xếp hạng trong dữ liệu nguồn}
\label{tab:nlp_rating_distribution}
\begin{tabular}{lr}
\toprule
Đánh giá & Số lượng đánh giá \\
\midrule
1 & 12,326 \\
2 & 1,873 \\
3 & 3,271 \\
4 & 6,692 \\
5 & 128,248 \\
\bottomrule
\end{tabular}
\end{table}

Nhãn tình cảm mạnh mẽ bao gồm xếp hạng 1, 2, 4 và 5. Điều này mang lại 149.139 hàng. Lớp phủ định có 14.199 hàng và lớp tích cực có 134.940 hàng. Tỷ lệ lớp là khoảng 9,5 đến 1. Sự mất cân bằng này rất quan trọng đối với các lựa chọn mô hình NLP.

\subsection{Tiếng ồn ngôn ngữ được quan sát}
Ngôn ngữ đánh giá chứa các dạng nhiễu NLP phổ biến. Tập dữ liệu bao gồm các dạng ngắn như gr8, thx, idk và các cụm từ không chính thức. Nó bao gồm nhiều thái cực khác nhau trong một câu, chẳng hạn như giao hàng tốt nhưng trễ. Nó bao gồm các cụm từ phủ định rõ ràng như không tệ và không tốt.

Tập dữ liệu cũng có các từ dành riêng cho tên miền dành cho thẻ quà tặng, đổi mã, kích hoạt, giao hàng, hỗ trợ và hoàn tiền. Những mô hình này thúc đẩy cả việc lập bản đồ từ vựng và học tập dựa trên dữ liệu.

\subsection{Dữ liệu nhãn phát hành}
Để trích xuất vấn đề, tập dữ liệu được gắn nhãn thủ công có 19.000 hàng. Bảng \ref{tab:nlp_issue_frequency} hiển thị tần số nhãn.

\begin{table}[H]
\centering
\caption{Tần suất nhãn phát hành trong bộ chú thích NLP thủ công}
\label{tab:nlp_issue_frequency}
\begin{tabular}{lrr}
\toprule
Nhãn & Đếm & Chia sẻ \\
\midrule
delivery\_shipping & 972 & 0.0512 \\
redemption\_activation & 284 & 0.0149 \\
product\_quality & 1,251 & 0.0658 \\
customer\_service & 182 & 0.0096 \\
refund\_return & 93 & 0.0049 \\
khả năng sử dụng & 873 & 0.0459 \\
value\_price & 558 & 0.0294 \\
fraud\_scam & 144 & 0.0076 \\
khác & 15,960 & 0.8400 \\
\bottomrule
\end{tabular}
\end{table}

Nhãn trung bình cho mỗi đánh giá là 1,069. Điều này hỗ trợ lập mô hình nhiều nhãn vì nhiều hàng không phải là văn bản một khái niệm.

\section{Thiết kế xử lý văn bản}
\subsection{Chuẩn hóa}
Quy trình chuẩn hóa cốt lõi sử dụng chuyển đổi chữ thường, xóa URL, dọn dẹp dấu câu và chuẩn hóa khoảng trắng. Điều này làm giảm sự thưa thớt của mã thông báo trong khi vẫn giữ được các từ cảm xúc cốt lõi.

Chuẩn hóa chữ viết tắt tùy chọn ánh xạ các dạng ngắn phổ biến thành dạng đầy đủ. Điều này bao gồm các ánh xạ như gr8 thành tuyệt vời, thx thành cảm ơn, u tới bạn và idk tới tôi không biết.

\subsection{Xử lý phủ định}
Phủ định là một thách thức quan trọng của NLP trong các bài đánh giá ngắn. Dự án sử dụng phương pháp gắn thẻ phủ định. Sau tín hiệu phủ định, một số mã thông báo tiếp theo sẽ nhận được tiền tố NOT\_. Kích thước cửa sổ là ba mã thông báo theo mặc định \cite{wiegand2010}.

Phương pháp này đơn giản nhưng hữu ích cho các mô hình thưa thớt. Nó giúp tách biệt tốt với NOT\_good và xấu với NOT\_bad trong không gian tính năng.

\subsection{Xử lý độ tương phản}
Đường dẫn cũng mô hình hóa độ tương phản diễn ngôn đơn giản. Khi các dấu hiệu như but, tuy nhiên, mặc dù, tuy nhiên, và chưa xuất hiện, văn bản có thể được chia thành các mệnh đề trái và phải. Sau đó, các tính năng mệnh đề riêng biệt và cờ đánh dấu sẽ được thêm vào.

Thiết kế này nắm bắt một mẫu đánh giá chung trong đó mệnh đề thứ hai mang ý kiến cuối cùng của người dùng.

\subsection{Mô hình cấp độ nhân vật}
Các đặc điểm ký tự từ 3 đến 5 gam được sử dụng làm sự kết hợp tùy chọn với các đặc điểm của từ. Điều này giúp khắc phục lỗi đánh máy, tiếng lóng ngắn và các biến thể chính tả nhỏ. Trong các tác vụ NLP trên văn bản người dùng, điều này thường cải thiện độ mạnh mẽ \cite{joulin2017}.

\section{Kỹ thuật tính năng cho NLP}
\subsection{Biểu diễn thưa thớt cổ điển}
Biểu diễn chính là TF-IDF với từ n-gram. Đường cơ sở sử dụng 1 đến 2 gam. Các biến thể khác kiểm tra từ 1 đến 3 gam, thẻ phủ định, mệnh đề tương phản, gam ký tự và đặc điểm từ vựng \cite{manning2008ir,wang2012}.

Bảng \ref{tab:nlp_variant_design} tóm tắt các biến thể ngữ cảnh đã được thử nghiệm.

\begin{table}[H]
\centering
\caption{Các biến thể của tính năng ngữ cảnh NLP}
\label{tab:nlp_variant_design}
\begin{tabular}{>{\raggedright\arraybackslash}p{1.2cm} >{\raggedright\arraybackslash}p{12cm}}
\toprule
Khác nhau & Thiết kế tính năng \\
\midrule
V0 & Từ cơ sở TF-IDF có phạm vi n-gram từ 1 đến 2 \\
V1 & Từ TF-IDF có phạm vi n-gram từ 1 đến 3 \\
V2 & Từ 1 đến 2 cộng với gắn thẻ phủ định \\
V3 & Mệnh đề từ 1 đến 2 cộng được phân tách bằng dấu tương phản \\
V4 & Từ 1 đến 2 cộng với gắn thẻ phủ định cộng với phân chia mệnh đề \\
V5 & Từ 1 đến 2 cộng với ký tự 3 đến 5 gram \\
V6 & Từ 1 đến 2 cộng với ký tự 3 đến 5 gram cộng với gắn thẻ phủ định \\
V7 & Từ 1 đến 2 cộng với gắn thẻ phủ định cộng với khối từ vựng cảm tính \\
\bottomrule
\end{tabular}
\end{table}

\subsection{Khối tính năng Lexicon}
Dự án bao gồm khối từ vựng cảm tính dựa trên AFINN trong một biến thể. Nó tính điểm thô, điểm nhận biết phủ định, điểm nhận biết cường độ và điểm có trọng số tương phản. Điều này thêm từ vựng rõ ràng trước biểu diễn thưa thớt \cite{nielsen2011}.

Phần này hữu ích cho việc học khóa học NLP vì nó kết nối các phương pháp từ vựng cổ điển với việc học có giám sát hiện đại trong một quy trình.

\section{Phương pháp mô hình hóa}
\subsection{Bộ phân loại tình cảm}
Công cụ phân loại tình cảm chính là Hồi quy logistic. Lựa chọn tính năng Chi bình phương được áp dụng với các giá trị K cao nhất từ 1.000 đến 10.000. Cài đặt trọng lượng lớp không bao gồm các biến thể không có, cân bằng và có trọng số \cite{forman2003}.

Các mô hình cảm tính bổ sung bao gồm Hồi quy logistic L1 cho hành vi nhúng thưa thớt, Cây quyết định và Rừng ngẫu nhiên.

\subsection{Mô hình trích xuất vấn đề}
Trích xuất vấn đề sử dụng đào tạo một so với phần còn lại trên chín nhãn. Mô hình mặc định là Hồi quy logistic với trọng số lớp cân bằng. Đường cơ sở LinearSVM cũng có sẵn \cite{rifkin2004,zhang2014}.

Việc điều chỉnh ngưỡng được thực hiện trên mỗi nhãn trên dữ liệu xác thực. Điều này rất quan trọng trong NLP nhiều nhãn vì một ngưỡng chung thường không đủ cho tất cả các nhãn.

\subsection{Máy biến áp mở rộng}
Phần mở rộng máy biến áp tinh chỉnh bộ phân loại trình tự với tổn thất theo trọng số lớp. Nó sử dụng văn bản rõ ràng và đánh giá với cùng một khung ngưỡng và độ không chắc chắn để so sánh trực tiếp \cite{devlin2019}.

Tiện ích mở rộng này cung cấp ngữ nghĩa theo ngữ cảnh sâu hơn và giúp phân tích các cụm từ ngắn khó hiểu trong đó các mô hình thưa thớt có thể thất bại.

\section{Thiết lập thử nghiệm}
\subsection{Phân chia và giao thức}
Tất cả các thử nghiệm đều sử dụng hạt giống 42. Các nhãn tình cảm mạnh mẽ được phân chia theo tỷ lệ đào tạo, xác thực và kiểm tra 70, 10 và 20. Đánh giá ba sao được giữ để phân tích độ không chắc chắn.

Việc trích xuất vấn đề sử dụng 19.000 hàng được gắn nhãn với kích thước phân chia là 13.300 cho tập huấn, 1.900 cho xác thực và 3.800 cho kiểm tra trong các lần chạy đã lưu.

\subsection{Số liệu chính}
Đối với cảm tính, các số liệu chính là độ chính xác\_0, thu hồi\_0, F1\_0 và F2\_0, cùng với độ chính xác hoàn toàn và F1. Các số liệu chọn lọc trong điều kiện không chắc chắn bao gồm phạm vi bao phủ, tỷ lệ không chắc chắn, độ chính xác có chọn lọc\_0, thu hồi có chọn lọc\_0, F2\_0 chọn lọc và tỷ lệ âm tính giả trên các mẫu được bao phủ \cite{saito2015}.

Để trích xuất vấn đề, các số liệu chính là micro-F1, macro-F1, độ chính xác của tập hợp con, tổn thất hamming và độ chính xác trên mỗi nhãn, thu hồi và F1.

\subsection{Quy tắc bất định}
Chính sách không chắc chắn sử dụng ngưỡng thấp và ngưỡng cao đối với xác suất dương. Nó có thể trở lại không chắc chắn ở phạm vi giữa. Nó cũng trả về sự không chắc chắn đối với các trường hợp văn bản quá ngắn hoặc thưa thớt trong quy trình cổ điển \cite{geifman2017}.

Quy tắc này làm cho đầu ra NLP an toàn hơn cho quá trình sản xuất vì văn bản có ít thông tin không bị ép vào một lớp sai cứng.

\section{Gói nâng cấp giáo trình không có dữ liệu mới}
Dự án hiện bao gồm một gói bổ sung giúp cải thiện sự liên kết với các chủ đề khóa học NLP mà không cần thêm bất kỳ dữ liệu được gắn nhãn mới nào. Gói này sử dụng lại cùng một chuỗi phân tách, xác thực và kiểm tra, đồng thời giữ nguyên chính sách dọn dẹp văn bản.

Khối nâng cấp đầu tiên là bộ điểm chuẩn cổ điển mạnh mẽ. Nó bao gồm Hồi quy logistic cân bằng, Bayes Naive đa thức, Bayes Naive bổ sung, Perceptron, LinearSVM, SGD với tính năng mất log, Hồi quy logistic ngữ nghĩa trên vectơ SVD và MLP chuyển tiếp nguồn cấp dữ liệu trên cùng vectơ SVD. Kích thước đoàn tàu được giới hạn ở 35.000 hàng cho khối này để duy trì thời gian chạy thực tế trong khi vẫn duy trì tỷ lệ lớp.

Khối nâng cấp thứ hai là khối mô hình ngôn ngữ. Nó đào tạo các mô hình ngôn ngữ unigram và bigram với khả năng làm mịn add-k và báo cáo sự phức tạp trong quá trình phân tách xác thực và kiểm tra. Điều này đưa ra bằng chứng trực tiếp về chủ đề mô hình ngôn ngữ N-gram trong khóa học \cite{jurafsky2026,manning2008ir}.

Khối nâng cấp thứ ba xây dựng ma trận bao phủ chủ đề từ các tạo phẩm của dự án. Ma trận sử dụng điểm số chủ đề rõ ràng và báo cáo điểm tổng thể về nội dung giáo trình để thảo luận bào chữa.

\section{Kết quả cho Mô hình tình cảm NLP cổ điển}
\subsection{Lựa chọn đường cơ sở bởi K}
Quét đường cơ sở trên V0 hiển thị cài đặt tốt nhất cho mỗi K theo quy tắc phủ định đầu tiên. Bảng \ref{tab:nlp_chi2_best_per_k} báo cáo các hàng đã chọn từ quá trình xác thực.

\begin{table}[H]
\centering
\caption{Kết quả tình cảm cơ bản tốt nhất trên mỗi K khi xác thực}
\label{tab:nlp_chi2_best_per_k}
\begin{tabular}{lcccccc}
\toprule
K & hạng cân & Precision\_0 & Recall\_0 & F2\_0 & Sự chính xác & F1 \\
\midrule
1,000 & w10 & 0.555 & 0.910 & 0.806 & 0.922 & 0.955 \\
2,000 & w10 & 0.584 & 0.920 & 0.825 & 0.930 & 0.960 \\
5,000 & cân bằng & 0.606 & 0.935 & 0.843 & 0.936 & 0.963 \\
10,000 & cân bằng & 0.621 & 0.936 & 0.850 & 0.940 & 0.966 \\
\bottomrule
\end{tabular}
\end{table}

\subsection{So sánh biến thể}
Sau khi quét toàn bộ biến thể, cài đặt cổ điển được chọn là V6 với K bằng 10.000 và trọng lượng hạng w10. Bảng \ref{tab:nlp_variant_test} hiển thị so sánh thử nghiệm giữa các biến thể.

\begin{table}[H]
\centering
\caption{Các biến thể tình cảm NLP cổ điển trên dữ liệu thử nghiệm}
\label{tab:nlp_variant_test}
\begin{tabular}{lcccccc}
\toprule
Khác nhau & K & Cân nặng & Precision\_0 & Recall\_0 & F2\_0 & Bảo hiểm \\
\midrule
V0 & 10,000 & cân bằng & 0.638 & 0.937 & 0.856 & 0.870 \\
V1 & 10,000 & cân bằng & 0.643 & 0.932 & 0.855 & 0.871 \\
V2 & 5,000 & w10 & 0.647 & 0.926 & 0.852 & 0.874 \\
V3 & 10,000 & cân bằng & 0.626 & 0.923 & 0.843 & 0.868 \\
V4 & 10,000 & cân bằng & 0.635 & 0.921 & 0.845 & 0.872 \\
V5 & 10,000 & cân bằng & 0.653 & 0.933 & 0.859 & 0.878 \\
V6 & 10,000 & w10 & 0.677 & 0.932 & 0.867 & 0.885 \\
V7 & 5,000 & w10 & 0.657 & 0.927 & 0.857 & 0.880 \\
\bottomrule
\end{tabular}
\end{table}

V6 mang lại sự cân bằng tốt cho tình cảm NLP trong dự án này. Nó duy trì khả năng thu hồi âm bản cao và cải thiện độ chính xác âm so với nhiều biến thể khác.

\subsection{Ngưỡng và hành vi không chắc chắn}
Đối với mẫu cổ điển V6 đã chọn, ngưỡng mặc định 0,40 và 0,60 mang lại phạm vi bảo hiểm 0,885 và thu hồi có chọn lọc\_0 0,956. Bảng \ref{tab:nlp_classic_threshold} hiển thị mức quét ngưỡng.

\begin{table}[H]
\centering
\caption{Quét ngưỡng cho mô hình tình cảm NLP cổ điển đã chọn}
\label{tab:nlp_classic_threshold}
\begin{tabular}{lcccc}
\toprule
Cặp ngưỡng & Bảo hiểm & Selective precision\_0 & Selective recall\_0 & Selective F2\_0 \\
\midrule
0.40 / 0.60 & 0.885 & 0.730 & 0.956 & 0.900 \\
0.45 / 0.60 & 0.891 & 0.708 & 0.956 & 0.894 \\
0.50 / 0.65 & 0.886 & 0.681 & 0.963 & 0.889 \\
0.55 / 0.70 & 0.881 & 0.646 & 0.970 & 0.882 \\
\bottomrule
\end{tabular}
\end{table}

\section{Kết quả cho phần mở rộng Transformer NLP}
\subsection{Số liệu máy biến áp}
Kết quả tinh chỉnh máy biến áp được lưu trong đầu ra mở rộng NLP. Bảng \ref{tab:nlp_transformer_metrics} tóm tắt các giá trị xác nhận và kiểm tra.

\begin{table}[H]
\centering
\caption{Số liệu mở rộng Transformer NLP}
\label{tab:nlp_transformer_metrics}
\begin{tabular}{lcccccc}
\toprule
Tách ra & Sự chính xác & F1 & Precision\_0 & Recall\_0 & Bảo hiểm & Selective recall\_0 \\
\midrule
Xác thực & 0.968 & 0.982 & 0.821 & 0.854 & 0.908 & 0.862 \\
Bài kiểm tra & 0.967 & 0.982 & 0.822 & 0.833 & 0.910 & 0.846 \\
\bottomrule
\end{tabular}
\end{table}

So với động cơ V6 cổ điển, máy biến áp cho độ chính xác tổng thể cao hơn và độ chính xác âm cao hơn ở ngưỡng cơ sở. Động cơ V6 cổ điển vẫn giữ mức thu hồi âm thô cao hơn nhiều.

\subsection{Quét ngưỡng máy biến áp}
Bảng \ref{tab:nlp_transformer_threshold} thể hiện hành vi chọn lọc của máy biến áp khi quét ngưỡng.

\begin{table}[H]
\centering
\caption{Quét ngưỡng máy biến áp trên dữ liệu thử nghiệm}
\label{tab:nlp_transformer_threshold}
\begin{tabular}{lcccc}
\toprule
Cặp ngưỡng & Bảo hiểm & Selective precision\_0 & Selective recall\_0 & Selective F2\_0 \\
\midrule
0.40 / 0.60 & 0.910 & 0.831 & 0.846 & 0.843 \\
0.45 / 0.60 & 0.911 & 0.827 & 0.847 & 0.843 \\
0.50 / 0.65 & 0.912 & 0.821 & 0.850 & 0.844 \\
0.55 / 0.70 & 0.911 & 0.818 & 0.855 & 0.847 \\
\bottomrule
\end{tabular}
\end{table}

Đường cong máy biến áp ổn định, có phạm vi bao phủ gần 0,91 và mức tăng vừa phải trong thu hồi có chọn lọc khi cặp ngưỡng di chuyển đến 0,55 và 0,70.

\subsection{So sánh trường hợp cứng}
Kết quả trường hợp khó trong Bảng \ref{tab:nlp_hard_cases} so sánh các quyết định về đường cơ sở và máy biến áp cổ điển.

\begin{table}[H]
\centering
\caption{So sánh trường hợp khó giữa mô hình tình cảm cổ điển và biến đổi}
\label{tab:nlp_hard_cases}
\begin{tabular}{>{\raggedright\arraybackslash}p{4.1cm} >{\centering\arraybackslash}p{2.3cm} >{\centering\arraybackslash}p{2.3cm} >{\raggedright\arraybackslash}p{4.3cm}}
\toprule
Chữ & Quyết định cơ bản & Quyết định máy biến áp & Quan sát chính \\
\midrule
không tệ & Tiêu cực & Tích cực & máy biến áp đọc thành ngữ là tích cực \\
không tốt & Tiêu cực & Tiêu cực & cả hai mô hình đều đồng ý về tiêu cực \\
tốt nhưng giao hàng chậm & Tiêu cực & Tích cực & biến áp trọng lượng đầu tiên mệnh đề tích cực hơn \\
gr8 & không chắc chắn & không chắc chắn & cả hai đều trả về văn bản ngắn không chắc chắn \\
cảm ơn & không chắc chắn & không chắc chắn & cả hai đều trả về văn bản ngắn không chắc chắn \\
tôi không biết & không chắc chắn & không chắc chắn & cả hai đều trả về văn bản ngắn không chắc chắn \\
\bottomrule
\end{tabular}
\end{table}

Bảng này cho thấy rằng máy biến áp xử lý một số mẫu ngữ nghĩa khác nhau. Nó có thể cải thiện các trường hợp thành ngữ chẳng hạn như không tệ, nhưng nó cũng có thể chuyển văn bản khiếu nại hỗn hợp sang hướng tích cực.

\section{Kết quả trích xuất vấn đề NLP}
\subsection{Kết quả đa nhãn tổng thể}
Bảng \ref{tab:nlp_issue_overall} so sánh ba thiết lập trích xuất vấn đề.

\begin{table}[H]
\centering
\small
\setlength{\tabcolsep}{4pt}
\caption{Số liệu tổng thể về trích xuất vấn đề NLP}
\label{tab:nlp_issue_overall}
\begin{tabularx}{\textwidth}{
>{\raggedright\arraybackslash}X
>{\centering\arraybackslash}p{2.1cm}
>{\centering\arraybackslash}p{1.8cm}
>{\centering\arraybackslash}p{1.8cm}
>{\centering\arraybackslash}p{2.4cm}}
\toprule
Cài đặt & Val vi-F1 & Kiểm tra micro-F1 & Kiểm tra macro-F1 & Kiểm tra độ chính xác của tập hợp con \\
\midrule
Word TF-IDF cộng với OVR Logistic & 0.759 & 0.765 & 0.381 & 0.606 \\
Word TF-IDF cộng với OVR LinearSVM & 0.750 & 0.757 & 0.371 & 0.602 \\
Word cộng với Char TF-IDF cộng với OVR Logistic & 0.764 & 0.763 & 0.408 & 0.603 \\
\bottomrule
\end{tabularx}
\end{table}

Việc biểu diễn từ cộng với ký tự mang lại macro-F1 cao hơn, trong khi thiết lập chỉ từ giữ độ chính xác của micro-F1 và tập hợp con cao hơn một chút trong lần chạy này.

\subsection{Cải thiện trên mỗi nhãn}
Bảng \ref{tab:nlp_issue_per_label_delta} hiển thị sự thay đổi F1 của thử nghiệm trên mỗi nhãn từ từ này sang từ khác cộng với thiết lập ký tự.

\begin{table}[H]
\centering
\caption{Thay đổi F1 kiểm tra trên mỗi nhãn trong trích xuất vấn đề}
\label{tab:nlp_issue_per_label_delta}
\begin{tabular}{lccc}
\toprule
Nhãn & từ F1 & Từ F1 cộng với char & Đồng bằng \\
\midrule
delivery\_shipping & 0.475 & 0.462 & -0.013 \\
redemption\_activation & 0.267 & 0.304 & +0.037 \\
product\_quality & 0.308 & 0.301 & -0.007 \\
customer\_service & 0.314 & 0.185 & -0.128 \\
refund\_return & 0.145 & 0.333 & +0.188 \\
khả năng sử dụng & 0.462 & 0.476 & +0.014 \\
value\_price & 0.284 & 0.316 & +0.032 \\
fraud\_scam & 0.265 & 0.388 & +0.123 \\
khác & 0.911 & 0.910 & -0.001 \\
\bottomrule
\end{tabular}
\end{table}

Các tính năng ký tự giúp một số nhãn thưa thớt như hoàn tiền\_return và lừa đảo\_scam, nhưng có thể giảm F1 trên một số nhãn. Điều này cho thấy sự đánh đổi về việc thu hồi độ chính xác theo nhãn.

\subsection{Chiến lược ngưỡng trong cài đặt nhiều nhãn}
Trích xuất vấn đề sử dụng ngưỡng điều chỉnh cho mỗi nhãn. Cấu hình ngưỡng khác nhau giữa thiết lập chỉ từ và thiết lập từ+char vì phân phối điểm thay đổi theo nhãn.

Hành vi này được mong đợi trong các tác vụ nhiều nhãn NLP trong đó mỗi nhãn có tần số trước và cường độ tín hiệu văn bản khác nhau.

\section{Kết quả cho Gói nâng cấp giáo trình}
\subsection{Đường cơ sở cổ điển cho các thuật toán khóa học cốt lõi}
Bảng \ref{tab:nlp_syllabus_bench} báo cáo số liệu kiểm tra cho khối cơ sở được thêm vào. Khối này sử dụng cùng một cách phân chia dữ liệu và không thêm nhãn mới.

\begin{table}[H]
\centering
\caption{Kết quả kiểm tra các cơ sở giáo trình NLP bổ sung}
\label{tab:nlp_syllabus_bench}
\begin{tabular}{lccc}
\toprule
Người mẫu & Precision\_0 & Recall\_0 & F2\_0 \\
\midrule
Hồi quy logistic (cân bằng) & 0.699 & 0.906 & 0.855 \\
Vịnh Naive đa thức & 0.663 & 0.895 & 0.836 \\
Bổ sung Naive Bayes & 0.531 & 0.941 & 0.815 \\
Perceptron & 0.686 & 0.842 & 0.806 \\
SVM tuyến tính & 0.795 & 0.829 & 0.822 \\
Mất nhật ký SGD & 0.618 & 0.915 & 0.835 \\
Hồi quy logistic ngữ nghĩa SVD & 0.505 & 0.911 & 0.785 \\
FFNN MLP trên vectơ SVD & 0.796 & 0.719 & 0.733 \\
\bottomrule
\end{tabular}
\end{table}

Những kết quả này cho thấy sự khác biệt rõ ràng về hành vi. Complement Naive Bayes có mức thu hồi âm cao nhất là 0,941, nhưng độ chính xác thấp ở mức 0,531. Hồi quy logistic cân bằng mang lại F2\_0 tốt nhất ở mức 0,855 với độ chính xác cao hơn. LinearSVM có độ chính xác cao nhất\_0 ở mức 0,795 nhưng mức thu hồi\_0 thấp hơn logistic. Điều này hữu ích cho kỳ thi và phòng thủ vì nó cho thấy sự cân bằng giữa các mô hình từ các chủ đề khóa học khác nhau trong một tập dữ liệu.

\subsection{Hành vi chọn lọc trong các đường cơ sở bổ sung}
Đối với các mô hình xác suất trong khối được thêm vào, áp dụng cùng một dải không chắc chắn 0,40 và 0,60. Hồi quy logistic cân bằng cho phạm vi kiểm tra 0,883 và thu hồi có chọn lọc\_0 0,932. Bổ sung Naive Bayes cho phạm vi kiểm tra 0,900 và thu hồi có chọn lọc\_0 0,952, nhưng độ chính xác chọn lọc\_0 chỉ là 0,568. Điều này xác nhận rằng khối mới phù hợp với chính sách của dự án và vẫn giữ được sự kiểm soát đánh đổi minh bạch.

\subsection{Kết quả mô hình ngôn ngữ N-Gram}
Bảng \ref{tab:nlp_ngram_results} báo cáo mức độ phức tạp của các mô hình unigram và bigram với add-k bằng 1,0.

\begin{table}[H]
\centering
\caption{Sự bối rối của khối mô hình ngôn ngữ N-gram}
\label{tab:nlp_ngram_results}
\begin{tabular}{lcc}
\toprule
Người mẫu & Xác thực bối rối & Kiểm tra sự bối rối \\
\midrule
Unigram bổ sung-k & 375.91 & 377.45 \\
Phần bổ sung Bigram-k & 358.24 & 362.72 \\
\bottomrule
\end{tabular}
\end{table}

Bigram mang lại độ phức tạp thấp hơn unigram trên cả hai phần tách. Điều này được mong đợi vì bigram sử dụng ngữ cảnh cục bộ từ mã thông báo trước đó. Khoảng cách không lớn vì văn bản ôn tập ngắn và ồn ào nhưng ổn định và có ý nghĩa đối với chủ đề mô hình ngôn ngữ.

\subsection{Ma trận phù hợp với khóa học được cập nhật}
Bảng \ref{tab:nlp_course_fit} hiển thị ma trận phạm vi chủ đề được cập nhật sau khối nâng cấp mới.

\begin{table}[H]
\centering
\caption{Phạm vi chủ đề được cập nhật sau khi nâng cấp không có dữ liệu mới}
\label{tab:nlp_course_fit}
\begin{tabular}{lcc}
\toprule
Nhóm chủ đề & Phần trăm bảo hiểm & Trạng thái \\
\midrule
Xác suất và đánh giá & 100 & mạnh \\
Mô hình ngôn ngữ N-gram & 100 & mạnh \\
Bayes ngây thơ và Perceptron & 100 & mạnh \\
Mô hình logistic và tuyến tính & 100 & mạnh \\
Mạng nơ-ron chuyển tiếp nguồn cấp dữ liệu & 100 & mạnh \\
Ngữ nghĩa và nhúng vectơ & 80 & vừa mạnh \\
Khái niệm cơ bản về máy biến áp và LLM & 100 & mạnh \\
Thực hành Pytorch & 60 & trung bình \\
RNN và LSTM & 20 & yếu đuối \\
Mô hình ngôn ngữ đeo mặt nạ & 20 & yếu đuối \\
Ứng dụng LLM & 30 & yếu đuối \\
\bottomrule
\end{tabular}
\end{table}

Điểm bảo hiểm tổng thể ước tính là 75,7\%. Con số này cao hơn nhiều so với trạng thái trước đó, nơi các chủ đề phân loại cốt lõi rất mạnh nhưng nhiều chủ đề bài kiểm tra vẫn còn thiếu trong bộ tạo phẩm.

\section{Thảo luận về máy biến áp cổ điển và máy biến áp}
NLP thưa thớt cổ điển và NLP biến áp cung cấp những điểm mạnh khác nhau. Động cơ V6 cổ điển duy trì khả năng thu hồi tiêu cực cao và thu hồi có chọn lọc mạnh mẽ trong bối cảnh dự án này. Máy biến áp cho độ chính xác toàn cầu cao hơn và độ chính xác âm mạnh hơn.

Đối với việc phân loại hỗ trợ cần nắm bắt càng nhiều tiêu cực càng tốt, động cơ V6 cổ điển sẽ phù hợp hơn với chính sách mục tiêu. Để xử lý sắc thái ngữ nghĩa và thành ngữ, biến áp có thể hữu ích và có thể được kết hợp với điều khiển dựa trên quy tắc.

Một con đường lai thực tế là có thể. Một lựa chọn là giữ mô hình cổ điển làm mô hình phân loại chính và sử dụng máy biến áp làm điểm ghi điểm phụ trong các trường hợp khó, không chắc chắn.

\section{Chế độ phân tích lỗi và lỗi}
Chế độ thất bại đầu tiên là sự phân cực hỗn hợp với ngôn ngữ tương phản. Các câu bắt đầu tích cực và kết thúc bằng lời phàn nàn có thể gây nhầm lẫn cho mô hình này hay mô hình khác tùy thuộc vào trọng tâm tính năng.

Chế độ lỗi thứ hai là văn bản rất ngắn với tín hiệu thấp. Dự án xử lý vấn đề này với lý do đầu ra và dự phòng không chắc chắn.

Kiểu lỗi thứ ba là vấn đề hiếm gặp về nhãn khi trích xuất nhiều nhãn. Các nhãn như hoàn tiền\_return và lừa đảo\_scam vẫn có ít sự hỗ trợ và có thể hiển thị điểm không ổn định nếu không có hỗ trợ ký tự hoặc nhiều dữ liệu chú thích hơn.

\section{Giá trị thực tiễn cho khóa học NLP}
Dự án này hiện cho thấy một tập hợp ý tưởng NLP thậm chí còn rộng hơn trong một hệ thống. Nó bao gồm chuẩn hóa văn bản, ánh xạ từ vựng, mã hóa thưa thớt n-gram, xử lý phạm vi phủ định, tín hiệu cấp mệnh đề, phân loại văn bản nhiều nhãn, hiệu chỉnh ngưỡng và tinh chỉnh biến áp. Giờ đây, nó cũng bao gồm các tạo phẩm mô hình ngôn ngữ Naive Bayes, Perceptron và N-gram rõ ràng, cùng với đường cơ sở mạng chuyển tiếp nguồn cấp dữ liệu trên các vectơ SVD ngữ nghĩa.

Dự án vẫn mang tính thiết thực cho công việc khóa học. Nó đã lưu các tạo phẩm, số liệu có thể tái tạo, quy trình làm việc chú thích và các tập lệnh suy luận có thể triển khai. Nó cũng có một ma trận bao phủ chủ đề có thể được hiển thị trực tiếp khi bào chữa.

Đối với việc học trong khóa học, phiên bản cập nhật cung cấp bằng chứng mạnh mẽ hơn về sự cân bằng giữa các phương pháp tiếp cận NLP cổ điển và sâu sắc theo một mục tiêu kinh doanh và một tập dữ liệu cố định.

\section{Hạn chế và công việc trong tương lai}
Hạn chế chính vẫn là sự mất cân bằng dữ liệu, đặc biệt là về tâm lý tiêu cực và các nhãn vấn đề hiếm gặp. Điều này ảnh hưởng đến độ chính xác và độ ổn định trong một số mô hình.

Một hạn chế khác là độ sâu của chủ đề. Dự án hiện bao gồm nhiều chủ đề khóa học hơn, nhưng RNN, nhiệm vụ mô hình ngôn ngữ mặt nạ và ứng dụng LLM vẫn còn nhẹ so với phần còn lại.

Công việc sau này vẫn có thể tránh được dữ liệu mới lúc đầu. Nó có thể thêm một đường cơ sở RNN nhỏ trên cùng một phần tách, thêm một đầu dò mô hình ngôn ngữ được che giấu với cùng một kho văn bản và thêm một bản demo ứng dụng LLM dựa trên lời nhắc hoặc truy xuất đơn giản. Điều này sẽ cải thiện mức độ bao phủ chủ đề trong khi vẫn giữ nguyên tập dữ liệu.

\section{Phần kết luận}
Dự án NLP này xây dựng một quy trình hiểu biết đánh giá hoàn chỉnh từ phân loại tình cảm đến trích xuất vấn đề và mở rộng máy biến áp. Mô hình thưa thớt nhận biết ngữ cảnh cổ điển mang lại hiệu suất tiêu cực đầu tiên mạnh mẽ. Tiện ích mở rộng biến áp mang lại độ chính xác tổng thể tốt hơn và hành vi ngữ nghĩa hữu ích trong một số cụm từ khó. Trích xuất vấn đề nhiều nhãn cung cấp tín hiệu hành động chi tiết từ văn bản đánh giá. Khối giáo trình được nâng cấp bổ sung thêm Naive Bayes, Perceptron, các đường cơ sở ngữ nghĩa chuyển tiếp nguồn cấp dữ liệu và đánh giá mô hình ngôn ngữ N-gram mà không cần thêm dữ liệu. Với bản nâng cấp này, dự án mang lại sự liên kết chặt chẽ hơn với các chủ đề khóa học và vẫn có thể tái sử dụng cho mục đích quốc phòng và thực tế.

\clearpage
\section*{Tài liệu tham khảo}
\addcontentsline{toc}{section}{Tài liệu tham khảo}
\begin{thebibliography}{99}

\bibitem{jurafsky2026}
Daniel Jurafsky and James H. Martin.
\textit{Speech and Language Processing} (3rd ed. draft).
Online manuscript, January 2026.
\newline\texttt{https://web.stanford.edu/~jurafsky/slp3/}

\bibitem{manning2008ir}
Christopher D. Manning, Prabhakar Raghavan, and Hinrich Schutze.
\textit{Introduction to Information Retrieval}.
Cambridge University Press, 2008.
\newline\texttt{https://nlp.stanford.edu/IR-book/}

\bibitem{turney2002}
Peter D. Turney.
Thumbs up or thumbs down? Semantic orientation applied to unsupervised classification of reviews.
In \textit{Proceedings of ACL}, 2002.
\newline\texttt{https://aclanthology.org/P02-1053/}

\bibitem{wang2012}
Sida I. Wang and Christopher D. Manning.
Baselines and Bigrams: Simple, Good Sentiment and Topic Classification.
In \textit{Proceedings of ACL}, 2012.
\newline\texttt{https://aclanthology.org/P12-2018/}

\bibitem{wiegand2010}
Michael Wiegand, Alexandra Balahur, Benjamin Roth, Dietrich Klakow, and Andres Montoyo.
A Survey on the Role of Negation in Sentiment Analysis.
In \textit{Proceedings of the Workshop on Negation and Speculation in Natural Language Processing}, 2010.
\newline\texttt{https://aclanthology.org/W10-3111/}

\bibitem{joulin2017}
Armand Joulin, Edouard Grave, Piotr Bojanowski, and Tomas Mikolov.
Bag of Tricks for Efficient Text Classification.
In \textit{Proceedings of EACL}, 2017.
\newline\texttt{https://aclanthology.org/E17-2068/}

\bibitem{nielsen2011}
Finn Arup Nielsen.
A new ANEW: Evaluation of a word list for sentiment analysis in microblogs.
\textit{arXiv preprint arXiv:1103.2903}, 2011.
\newline\texttt{https://arxiv.org/abs/1103.2903}

\bibitem{forman2003}
George Forman.
An extensive empirical study of feature selection metrics for text classification.
\textit{Journal of Machine Learning Research}, 3:1289--1305, 2003.
\newline\texttt{https://www.jmlr.org/papers/v3/forman03a.html}

\bibitem{rifkin2004}
Ryan Rifkin and Aldebaro Klautau.
In defense of one-vs-all classification.
\textit{Journal of Machine Learning Research}, 5:101--141, 2004.
\newline\texttt{https://www.jmlr.org/papers/v5/rifkin04a.html}

\bibitem{tsoumakas2011}
Grigorios Tsoumakas, Ioannis Katakis, and Ioannis Vlahavas.
Random k-Labelsets for multilabel classification.
\textit{IEEE Transactions on Knowledge and Data Engineering}, 23(7):1079--1089, 2011.
\newline\texttt{https://doi.org/10.1109/TKDE.2010.164}

\bibitem{zhang2014}
Min-Ling Zhang and Zhi-Hua Zhou.
A review on multi-label learning algorithms.
\textit{IEEE Transactions on Knowledge and Data Engineering}, 26(8):1819--1837, 2014.
\newline\texttt{https://doi.org/10.1109/TKDE.2013.39}

\bibitem{vaswani2017}
Ashish Vaswani, Noam Shazeer, Niki Parmar, Jakob Uszkoreit, Llion Jones, Aidan N. Gomez, Lukasz Kaiser, and Illia Polosukhin.
Attention Is All You Need.
In \textit{Advances in Neural Information Processing Systems 30}, 2017.
\newline\texttt{https://arxiv.org/abs/1706.03762}

\bibitem{devlin2019}
Jacob Devlin, Ming-Wei Chang, Kenton Lee, and Kristina Toutanova.
BERT: Pre-training of Deep Bidirectional Transformers for Language Understanding.
In \textit{Proceedings of NAACL-HLT}, 2019.
\newline\texttt{https://aclanthology.org/N19-1423/}

\bibitem{saito2015}
Takaya Saito and Marc Rehmsmeier.
The precision-recall plot is more informative than the ROC plot when evaluating binary classifiers on imbalanced datasets.
\textit{PLOS ONE}, 10(3):e0118432, 2015.
\newline\texttt{https://doi.org/10.1371/journal.pone.0118432}

\bibitem{geifman2017}
Yonatan Geifman and Ran El-Yaniv.
Selective Classification for Deep Neural Networks.
\textit{arXiv preprint arXiv:1705.08500}, 2017.
\newline\texttt{https://arxiv.org/abs/1705.08500}

\end{thebibliography}

\clearpage
\section*{Phụ lục A: Số liệu tình cảm}
\addcontentsline{toc}{section}{Phụ lục A: Số liệu tình cảm}

\begin{figure}[H]
\centering
\includegraphics[width=0.82\textwidth]{\detokenize{results/dm2_steps/01_rating_distribution.png}}
\caption{Phân phối xếp hạng trong dữ liệu nguồn}
\end{figure}

\begin{figure}[H]
\centering
\includegraphics[width=0.82\textwidth]{\detokenize{results/dm2_steps/06_f1_vs_k.png}}
\caption{Xác thực F1 bằng Chi-square top K}
\end{figure}

\begin{figure}[H]
\centering
\includegraphics[width=0.82\textwidth]{\detokenize{results/dm2_steps/06_recall0_vs_k.png}}
\caption{Thu hồi xác thực cho lớp âm theo top K}
\end{figure}

\begin{figure}[H]
\centering
\includegraphics[width=0.82\textwidth]{\detokenize{results/dm2_steps/06b_recall0_by_variant.png}}
\caption{Kiểm tra việc thu hồi lớp phủ định bằng biến thể NLP cổ điển}
\end{figure}

\begin{figure}[H]
\centering
\includegraphics[width=0.82\textwidth]{\detokenize{results/dm2_steps/06b_precision0_by_variant.png}}
\caption{Kiểm tra độ chính xác cho lớp phủ định bằng biến thể NLP cổ điển}
\end{figure}

\begin{figure}[H]
\centering
\includegraphics[width=0.82\textwidth]{\detokenize{results/dm2_steps/09_prob_hist_test_strong.png}}
\caption{Biểu đồ xác suất mô hình cổ điển trên tập kiểm tra mạnh}
\end{figure}

\begin{figure}[H]
\centering
\includegraphics[width=0.82\textwidth]{\detokenize{results/dm2_steps/10_threshold_tradeoff.png}}
\caption{Sự đánh đổi ngưỡng mô hình cổ điển cho mức độ phù hợp và thu hồi}
\end{figure}

\begin{figure}[H]
\centering
\includegraphics[width=0.82\textwidth]{\detokenize{results/nlp_ext/nlp_confusion_matrix.png}}
\caption{Ma trận nhầm lẫn máy biến áp trên bộ thử nghiệm}
\end{figure}

\begin{figure}[H]
\centering
\includegraphics[width=0.82\textwidth]{\detokenize{results/nlp_ext/nlp_prob_hist_test.png}}
\caption{Biểu đồ xác suất máy biến áp trên tập thử nghiệm}
\end{figure}

\begin{figure}[H]
\centering
\includegraphics[width=0.82\textwidth]{\detokenize{results/nlp_ext/nlp_threshold_tradeoff.png}}
\caption{Sự đánh đổi ngưỡng máy biến áp để phủ sóng và thu hồi}
\end{figure}

\clearpage
\section*{Phụ lục B: Số liệu trích xuất vấn đề}
\addcontentsline{toc}{section}{Phụ lục B: Số liệu trích xuất vấn đề}

\begin{figure}[H]
\centering
\includegraphics[width=0.82\textwidth]{\detokenize{results/issue_steps/02_label_frequency.png}}
\caption{Tần suất phát hành nhãn trong bộ chú thích thủ công}
\end{figure}

\begin{figure}[H]
\centering
\includegraphics[width=0.82\textwidth]{\detokenize{results/issue_steps/02_per_label_f1.png}}
\caption{Mỗi nhãn F1 cho mô hình trích xuất vấn đề từ}
\end{figure}

\begin{figure}[H]
\centering
\includegraphics[width=0.82\textwidth]{\detokenize{results/issue_steps_char_demo/02_per_label_f1.png}}
\caption{Mỗi nhãn F1 cho mô hình trích xuất vấn đề từ cộng với char}
\end{figure}

\clearpage
\section*{Phụ lục C: Số liệu nâng cấp giáo trình}
\addcontentsline{toc}{section}{Phụ lục C: Số liệu nâng cấp giáo trình}

\begin{figure}[H]
\centering
\includegraphics[width=0.82\textwidth]{\detokenize{results/nlp_ext/syllabus_upgrade/nlp_syllabus_bench_recall0.png}}
\caption{Kiểm tra thu hồi\_0 trên các cơ sở giáo trình bổ sung}
\end{figure}

\begin{figure}[H]
\centering
\includegraphics[width=0.82\textwidth]{\detokenize{results/nlp_ext/syllabus_upgrade/nlp_syllabus_bench_f2.png}}
\caption{Kiểm tra F2\_0 trên các cơ sở giáo trình bổ sung}
\end{figure}

\end{document}
